\section{A vector-valued hashing lemma}\label{sec:hashing}

The following lemma turns dispersion among characters into dispersion among
table entries.  Its constants are independent of the dimension of the character
space and of the dimension of the vectors being hashed. In the application,
$q$ will index candidate labels and $z$ will index entries of a gate table.
The resulting hashed vector will collect the Fourier coefficients
contributing to one entry's influence, so its squared norm will be that influence. This is why
we need a vector-valued statement rather than a scalar estimate.

\begin{lemma}[Vector-valued hashing]\label{lem:hashing}
Let $H$ be a finite-dimensional vector space over $\F$, write
$\chi_q(u)=(-1)^{q(u)}$ for $q\in H^*$ and $u\in H$, let $\mathcal V$ be a
real Hilbert space, let $s\ge1$ be an integer, and let
\[
 c(q,z)\in\mathcal V,
 \qquad q\in H^*,\quad z\in\F^s.
\]
Set
\[
 w_q=\sum_{z\in\F^s}\norm{c(q,z)}^2,
 \qquad M=\sum_{q\in H^*}w_q.
\]
Choose $L,J_1,\ldots,J_s$ independently and uniformly in $H$, and define
\[
 \mathbf J(q)=(q(J_1),\ldots,q(J_s)),
 \qquad
 U(z)=\sum_{q\in H^*}\chi_q(L)c(q,z+\mathbf J(q)).
\]
Then
\begin{equation}\label{eq:hashing-total-mass}
 \E\sum_{z\in\F^s}\norm{U(z)}^2=M.
\end{equation}
For every $\zeta>0$, there are $\gamma>0$ and $s_0\in\mathbb N$, depending
only on $\zeta$, such that for every $s\ge s_0$,
\begin{equation}\label{eq:hashing-dispersion}
 \max_{q\in H^*}w_q\le\gamma M
 \quad\Longrightarrow\quad
 \E\max_{z\in\F^s}\norm{U(z)}^2\le\zeta M.
\end{equation}
\end{lemma}

\begin{proof}
All expectations over $H$ and $\F^s$ below use uniform probability measure;
the displayed sums defining $w_q$ and $M$ are unnormalized.  Conditional on
the $J_g$, orthogonality of the characters of $L$ gives
\[
 \E_L\norm{U(z)}^2
 =\sum_q\norm{c(q,z+\mathbf J(q))}^2.
\]
Summing over $z$ proves \eqref{eq:hashing-total-mass}.  If $M=0$, every vector
$c(q,z)$ is zero and \eqref{eq:hashing-dispersion} is immediate.  Henceforth
assume $M>0$.

Set $h=2^s$ and $\rho=\max_q w_q/M$. We will prove the quantitative
bound below for every $C,D>0$, and then choose its parameters.
The main estimate bounds the maximum squared norm by the square root of
the sum of fourth powers. After a Fourier change of variables, most terms
in that fourth moment force the same character $q$ in all four factors;
their total is controlled by the dispersion of the character masses.
Two truncations will control the remaining terms and the cost of this
change to the array:
\begin{equation}\label{eq:hashing-quantitative}
 \E\max_z\norm{U(z)}^2
 \le 2M\left[
 \sqrt{2\rho+\frac{2C^2}{D^2}+\frac{3D^4}{h}}
 +\left(\frac1{\sqrt C}+\frac1D\right)^2+\frac1h
 \right].
\end{equation}


\paragraph{Fourier representation.}
For $p\in\F^s$ and $u\in H$, define
\[
 \alpha_q(p)=\sum_{z\in\F^s}(-1)^{p\cdot z}c(q,z),
 \qquad
 A_p(u)=\sum_{q\in H^*}\chi_q(u)\alpha_q(p),
 \qquad
 S_p=L+\sum_{g=1}^s p_gJ_g.
\]
Fourier inversion in $z$ gives
\begin{equation}\label{eq:hashing-inversion}
 U(z)=\E_p(-1)^{p\cdot z}A_p(S_p).
\end{equation}
Writing $m_p=\E_u\norm{A_p(u)}^2$, Parseval's identity in the two groups gives
\begin{equation}\label{eq:hashing-parseval}
 m_p=\sum_q\norm{\alpha_q(p)}^2,
 \qquad \E_p m_p=M,
 \qquad \E_p\norm{\alpha_q(p)}^2=w_q.
\end{equation}
For distinct $p,p'$, the variables $S_p,S_{p'}$ are independent and uniform
in $H$: the coefficient vectors $(1,p),(1,p')$ are linearly independent
over $\F$.  Similarly, any three distinct $p_1,p_2,p_3$ give three
independent uniform variables $S_{p_1},S_{p_2},S_{p_3}$.  Indeed, any
nonempty subcollection of their coefficient vectors whose sum is zero
must have even cardinality, and hence would consist of two equal vectors.

\paragraph{Two truncations.}
We first discard indices $p$ with unusually large mean squared norm
$m_p$, and then discard unusually large values of $A_p(u)$ on the
remaining indices. The first cutoff makes the change to each individual
character coefficient uniformly small under the second cutoff. The
second gives the pointwise bound needed when the fourth-moment indices
coincide. We will control the discarded part using a first-moment bound
and a variance bound from the pairwise independence of the samples $S_p$.
Define
\[
 A'_p(u)=
 \begin{cases}
 A_p(u),&m_p\le CM\text{ and }\norm{A_p(u)}\le D\sqrt M,\\
 0,&\text{otherwise},
 \end{cases}
\]
and put
\[
 \alpha'_q(p)=\E_u\chi_q(u)A'_p(u),
 \qquad w'_q=\E_p\norm{\alpha'_q(p)}^2,
 \qquad U'(z)=\E_p(-1)^{p\cdot z}A'_p(S_p).
\]
Pointwise truncation does not increase the total squared norm, so Parseval
gives $\sum_qw'_q\le M$.  If $m_p\le CM$, then
\[
 \norm{\alpha_q(p)-\alpha'_q(p)}
 \le\E_u\norm{A_p(u)}
       \one_{\{\norm{A_p(u)}>D\sqrt M\}}
 \le\frac{m_p}{D\sqrt M}
 \le\frac{C\sqrt M}{D}.
\]
If $m_p>CM$, then $\alpha'_q(p)=0$.  Consequently,
\begin{equation}\label{eq:hashing-truncated-mass}
 \sum_qw'_q\le M,
 \qquad
 w'_q\le 2w_q+\frac{2C^2M}{D^2}.
\end{equation}

Let
\[
 X_p=\norm{(A_p-A'_p)(S_p)},\qquad X=\E_pX_p.
\]
Equation~\eqref{eq:hashing-inversion} and the triangle inequality imply
$\max_z\norm{U(z)-U'(z)}\le X$.  The contribution to $\E X$ from the
indices with $m_p>CM$ is at most
\[
 \E_p\one_{\{m_p>CM\}}\sqrt{m_p}
 \le\frac{\E_p m_p}{\sqrt{CM}}=\sqrt{M/C}.
\]
The contribution from the remaining indices is at most
$\E_p m_p/(D\sqrt M)=\sqrt M/D$.  Moreover, the $X_p$ are pairwise
independent, by the pairwise independence of the $S_p$.  Since
$A_p-A'_p$ is either $A_p$ or zero at each point,
\[
 \operatorname{Var}(X)
 =h^{-2}\sum_p\operatorname{Var}(X_p)
 \le h^{-2}\sum_p\E_u\norm{A_p(u)-A'_p(u)}^2
 \le\frac Mh.
\]
We have proved
\begin{equation}\label{eq:hashing-truncation-error}
 \E\max_z\norm{U(z)-U'(z)}^2
 \le M\left[\left(\frac1{\sqrt C}+\frac1D\right)^2+\frac1h\right].
\end{equation}

\paragraph{The fourth moment.}
For brevity write $B_p=A'_p(S_p)$.  Expanding the fourth power of the norm
and summing the characters in $z$ gives the exact identity
\begin{align}
 \E\sum_z\norm{U'(z)}^4
 &=h^{-3}\sum_{p_1,p_2,p_3}
   \E\bigl[\ip{B_{p_1}}{B_{p_2}}
                  \ip{B_{p_3}}{B_{p_1+p_2+p_3}}\bigr]\notag\\
 &=\E_{p_1,p_2,p_3}\E
   \bigl[\ip{B_{p_1}}{B_{p_2}}\ip{B_{p_3}}{B_{p_4}}\bigr],
 \qquad p_4=p_1+p_2+p_3.
 \label{eq:hashing-fourth-expansion}
\end{align}
Here $p_1,p_2,p_3$ are independent uniform indices in the second line.
The factor $h^{-3}$ in the first line comes from four factors $h^{-1}$
in Fourier inversion and the sum of $h$ characters over $z$.

If any two of $p_1,p_2,p_3,p_4$ coincide, then some two of
$p_1,p_2,p_3$ coincide.  This event has probability at most $3/h$.
Since $\norm{B_p}\le D\sqrt M$, its contribution to
\eqref{eq:hashing-fourth-expansion} is at most $3D^4M^2/h$ in absolute value.

Now fix four distinct indices with $p_4=p_1+p_2+p_3$.  The first three
$S_{p_i}$ are independent uniform elements of $H$, and
$S_{p_4}=S_{p_1}+S_{p_2}+S_{p_3}$.  Expanding the four $A'_{p_i}$ in
characters and averaging over those three independent elements yields
\begin{equation}\label{eq:hashing-common-character}
 \E\bigl[\ip{B_{p_1}}{B_{p_2}}\ip{B_{p_3}}{B_{p_4}}\bigr]
 =\sum_q
   \ip{\alpha'_q(p_1)}{\alpha'_q(p_2)}
   \ip{\alpha'_q(p_3)}{\alpha'_q(p_4)}.
\end{equation}
In detail, characters indexed by $q_1,q_2,q_3,q_4$ produce the factor
\[
 \prod_{i=1}^3\chi_{q_i+q_4}(S_{p_i}),
\]
whose expectation is zero unless $q_1=q_2=q_3=q_4$.

The right side of \eqref{eq:hashing-common-character} need not be
nonnegative.  For an upper bound, replace each summand by
$\prod_{i=1}^4\norm{\alpha'_q(p_i)}$, which is nonnegative, and then
drop the restriction that the four indices be distinct.  For every fixed
$q$, Cauchy--Schwarz gives
\begin{align*}
 \E_{p_1,p_2,p_3}\prod_{i=1}^4\norm{\alpha'_q(p_i)}
 &\le
 \left(\E\norm{\alpha'_q(p_1)}^2
             \norm{\alpha'_q(p_2)}^2\right)^{1/2}
 \left(\E\norm{\alpha'_q(p_3)}^2
             \norm{\alpha'_q(p_4)}^2\right)^{1/2}\\
 &=(w'_q)^2.
\end{align*}
The last equality uses the unconditioned independent uniform law of each
pair $(p_1,p_2)$ and $(p_3,p_4)$; it does not assert independence after
conditioning on distinctness. Thus the nondegenerate fourth moment is
bounded by $\sum_q(w'_q)^2$: dispersion among characters now controls
concentration among table entries. Combining these estimates with
\eqref{eq:hashing-truncated-mass} proves
\begin{align*}
 \E\sum_z\norm{U'(z)}^4
 &\le \sum_q(w'_q)^2+\frac{3D^4M^2}{h}\\
 &\le M^2\left(2\rho+\frac{2C^2}{D^2}+\frac{3D^4}{h}\right).
\end{align*}
Therefore, using Cauchy--Schwarz once more,
\begin{equation}\label{eq:hashing-truncated-maximum}
 \E\max_z\norm{U'(z)}^2
 \le\left(\E\sum_z\norm{U'(z)}^4\right)^{1/2}
 \le M\sqrt{2\rho+\frac{2C^2}{D^2}+\frac{3D^4}{h}}.
\end{equation}
The inequality $(a+b)^2\le2a^2+2b^2$, together with
\eqref{eq:hashing-truncation-error} and
\eqref{eq:hashing-truncated-maximum}, proves
\eqref{eq:hashing-quantitative}.

Finally, set $\varepsilon=\min\{\zeta,1\}$ and choose
\[
 C=\frac{64}{\varepsilon},\qquad
 D=\frac{8C}{\varepsilon},\qquad
 \gamma=\frac{\varepsilon^2}{128},\qquad
 2^{s_0}\ge
 \max\left\{\frac{192D^4}{\varepsilon^2},\frac{16}{\varepsilon}\right\}.
\]
If $s\ge s_0$ and $\rho\le\gamma$, the square-root term in
\eqref{eq:hashing-quantitative} is at most $\varepsilon/4$.
Also $D\ge\sqrt C$, so the other two terms in its brackets are at most
$\varepsilon/16$ each.  The resulting bound is
$3\varepsilon M/4\le\zeta M$, establishing
\eqref{eq:hashing-dispersion} with constants depending only on $\zeta$.
\end{proof}
